\documentclass[../main.tex]{subfiles}
\begin{document}
%==========================================TIÊU ĐỀ TẬP========================================
\begin{table}[h]
  \begin{adjustwidth}{-1cm}{}
    \begin{tabular}{>{\centering\arraybackslash}p{6cm}|>{\raggedright\arraybackslash}p{12.5cm}}
      \multirow{5}{6cm}
      % Cột bên trái: ÉPISODE và số tập
      {\\[-40pt]
        \flaregothic\fontsize{20pt}{20pt}\selectfont \centering
        \color{eptype}ÉPISODE\color{black}\\
        \burbank\fontsize{72pt}{72pt}\selectfont
        \color{epnum}132\color{black}\\[6pt]
        \cafeteria\fontsize{20pt}{20pt}\selectfont\color{epnum}(10-11)\color{black}
        \\[18pt]
        % \flaregothic\fontsize{14pt}{14pt}\selectfont
        % \color{epattr}BIENVENUE, \uline{\color{Banpire} Mel} !
        % \flaregothic\fontsize{18pt}{18pt}\selectfont
        % \color{epattr}ÉPISODE PERDU
        \color{black}
      }
       &
      % Dòng thứ nhất: tên tập tiếng Nhật
      {\fotskip\fontsize{18pt}{18pt}\selectfont そんなもの\ruby{読}{よ}んじゃいけません！
      }                                     \\[6pt]
      % Dòng thứ hai: tên tập tiếng Anh
       & {\cafeteria\fontsize{18pt}{18pt}\selectfont Careful What You Read\dots}                             \\[4pt]
      % Dòng thứ ba: tên tập tiếng Việt
       & {\vnmsans\fontsize{18pt}{18pt}\selectfont Giang hồ Ayame}                                           \\[8pt]
      % Dòng thứ tư: Nhân vật xuất hiện
       & {\sen{\fontsize{14pt}{14pt}\selectfont Track used: Rhythm Party/Beat party (リズムのパーティー/リズミックパーティー)}}
      \\[6pt]
      % Dòng cuối cùng: Thời gian diễn ra của tập (thường sẽ là ngày phát hành)
       & {\continuum\fontsize{16pt}{16pt}\selectfont Ngày phát hành: 21/11/2021}                             \\[8pt]
       & (dựa theo bản dịch của \href{https://www.facebook.com/mamisubteam}{MamiSubs - Mami Fansub})         \\[18pt]
    \end{tabular}
  \end{adjustwidth}
\end{table}
%=========================================TRUYỆN CHÍNH========================================
\color{black}

Chiều muộn, vào một ngày nửa cuối tháng Mười Một, tiết trời bắt đầu trở lạnh, không khí ngoài trời nhuốm một màu vàng xám uể oải của hoàng hôn mùa đông. Văn phòng \hll\ cũng dần thưa người, hầu hết các nhân viên và thành viên đã về gần hết từ lúc xế chiều, chỉ còn lại vài bóng người lác đác ở lại để hoàn tất nốt vài công việc lặt vặt. Sự yên ắng kỳ lạ khiến cả tầng như chìm vào một cơn buồn ngủ tập thể.

Tại khu văn phòng chính, Subaru đang nằm dài trên ghế sofa với một cuốn truyện tranh trong tay. Đó là một trong những tuyển tập yêu thích của cô, xoay quanh đề tài xã hội đen và giới giang hồ, với những lời thoại đậm chất mafia Nhật Bản.

``Quạu à\dots'' nàng Tomboy vừa lật trang vừa ngẫm nghĩ, miệng lẩm bẩm như đang hòa mình vào vai một tên đàn em đầu gấu nào đó.

Bất chợt, khi ánh nhìn từ trang truyện lướt ra ngoài cửa sổ, Subaru bàng hoàng nhận ra trời đã nhá nhem từ khi nào. ``\textbf{Thôi chết, muộn giờ rồi!!}'' cô thốt lên, bật dậy theo phản xạ, vô tình hất cuốn truyện bay vút lên không trung.

*BỤP!*

Cuốn truyện ấy, như định mệnh an bài, bay một đường cong hoàn hảo rồi đập đúng trán Ayame, người đang say giấc trong tư thế nghiêng đầu trên ghế dài ở phía đối diện.

\img{10-11a}

``Cái gì đây\dots?'' nàng Quỷ sừng dụi mắt, nhăn mặt vì bị đánh thức bất ngờ. Cô cầm cuốn truyện lên, định ném lại cho Subaru thì đập vào mắt cô là hình ảnh một nữ thủ lĩnh yakuza đội mũ lưỡi trai và đeo kính đen ngầu lòi.

Hình ảnh đó, không hiểu sao, lập tức chạm đến một góc khuất sâu trong tâm hồn của Ayame, nơi ẩn chứa một tính cách có phần hiếu kì và ham học hỏi như một đứa trẻ mới lớn.

Thế là cô ngồi nguyên tại chỗ, nghiền ngẫm từng trang truyện, từng tập truyện cho đến tận khi văn phòng chỉ còn lại tiếng điều hòa chạy đêm.

\begin{center}
  \uline{Sáng hôm sau.}
\end{center}

Subaru bước vào phòng với vẻ mặt đầy năng lượng sau một đêm ngủ bù.

``Có gì mới nào?'' cô cất tiếng chào, không biết rằng mình sắp đối mặt với một thứ còn sốc hơn cả việc bị trễ hạn.

Ngay sau lưng là Kuon, cậu quản lý trung thành, đang kéo cái túi đựng tài liệu theo một cách uể oải. ``Hy vọng là hôm nay không có gì quá quái dị nữa\dots'' cậu thì thầm, như muốn xua đi điềm gở.

Nhưng hy vọng ấy đã tan thành mây khói ngay khi cả hai nhìn thấy Ayame đang ngồi gác chân lên bàn, diện trong bộ trang phục có phần ngược đời hơn thường ngày.

Diện trong bộ áo kimono quen thuộc mà mọi người thường thấy, nhưng nàng Quỷ sừng đã khoác lên mình một tấm áo hoàn toàn khác. Trên đầu cô là một chiếc mũ lưỡi trai đen tuyền, có đục sẵn hai lỗ để đôi sừng của cô chui qua, chính giữa là hình tam giác xanh được viền mạ vàng, xung quanh gắn hình ba bông hoa bốn cánh với ba màu vàng, đỏ và tím rực rỡ. Cặp kính râm kiểu thập niên 90 che gần nửa khuôn mặt, còn miệng thì đang ngậm một cành mầm lá nhỏ.

``Yo, Subaru, Kuon,'' ô cất giọng, khàn khàn và trầm, giống hệt mấy tay đàn anh trong giới giang hồ ở các bộ phim cổ lỗ sĩ.

\img{10-11b}

Nàng Vịt giật mình, suýt đánh rơi điện thoại trên tay: ``Bà\dots\ bà đang mặc cái gì thế, Ayame!?''

Cậu Tổng đứng bên cạnh, trán toát mồ hôi, không thể tin nổi vào mắt mình: ``Trông em như bước ra từ một manga giang hồ nào đó vậy\dots''

``Thì trang phục mới của ta là như vậy mà,'' nàng Oni đáp tỉnh queo, rồi lại cúi xuống ngồi xổm, làm như mình là đại ca đứng canh cổng bến cảng.

Subaru và Kuon quay sang nhìn nhau đầy hoài nghi, quay sang hội ý riêng.

``Bà ấy học đâu ra cái kiểu bố đời này vậy?'' nàng Vịt lẩm bẩm.

``Anh chịu,'' Kuon nhún vai, lắc đầu, mắt vẫn chớp chớp vì chưa hiểu nổi cách ăn mặc của Ayame. ``Tự dưng làm anh nhớ tới\dots\ con Rồng nào đó đã nghỉ việc từ mấy tháng trước rồi.''

``Ý anh là\dots\ \emph{em ấy} sao?''

``Kiểu vậy.''

\il

\begin{center}
  \uline{Nhà Hikawa.}
\end{center}

Coco và Aloe đang đứng trong căn phòng giặt đồ mới sửa, cùng nhau gập quần áo và giặt thêm mẻ mới.

Đột nhiên, nàng Rồng tóc cam hắt hơi một tiếng rõ to. ``\textbf{Hắt xì!!!}''

``Chị sao vậy, Coco?'' nàng Succubus quay sang hỏi.

Coco khịt mũi: ``Chị cũng không biết\dots\ Cảm giác như ai đó đang gọi tên mình vậy\dots''

``Hình như em nghe nói chị từng là kiểu `giang hồ hạng sang' hay cái gì đó thì phải?''

``Đúng là chị đã từng như vậy thật\dots\ \textbf{Hắt xì!} Nhưng mà chuyện đó đâu liên quan- \textbf{Hắt xì!}''

``\dots Hay là do Kuon-kun nhắc đến chi?''

``Chị không nghĩ Kuon-paisen sẽ- \textbf{Hắt xì!} -nhắc mấy cái này đâu\dots''

Câu chuyện cứ thế tiếp diễn, kèm thêm vài ba cú hắt xì nữa từ Coco trong vòng mười phút, khiến Aloe càng lúc càng nghi ngờ. ``\textit{Kể từ khi chị ấy chuyển đến sống ở đây, chưa bao giờ hắt xì nhiều đến vậy,}'' cô nghĩ thầm.

\il

\begin{center}
  \uline{Văn phòng.}
\end{center}

``Anh vẫn không nghĩ Ayame-chan học cái phong cách này từ Coco đâu\dots'' Kuon lắc đầu. ``Em ấy tốt nghiệp tầm năm tháng trước rồi mà, với lại anh không nghĩ ai lại biết nhà em ấy giờ ở đâu\dots''

``Vậy thì rốt cuộc là từ đâu ra nhỉ\dots'' Subaru ngẫm nghĩ, rồi vô tình liếc xuống cuốn truyện cô từng đọc hôm qua. Bìa truyện có một nữ thủ lĩnh yakuza đội mũ, đeo kính đen, trông giống bộ dạng mà nàng Quỷ sừng hiện tại.

``Khoan đã\dots\ Bà đọc cái này hả?'' cô hỏi, giơ cuốn truyện mà cô đọc dở lên.

Kuon tiếp lời, mắt đảo qua đảo lại giữa bìa truyện và nàng Quỷ sừng, gật đầu xác nhận: ``Nhìn vào bìa thì anh dám cá 99\% là vậy.''

``Chuẩn đét luôn!'' Ayame đứng dậy đầy khí thế. ``Đây là Kinh thánh của ta mà!''

Nàng Tomboy thở dài: ``Thế là học theo nguyên si không sót một chi tiết nào luôn\dots''

``Anh còn chẳng dám nghĩ em học được cái gì từ nó\dots'' cậu Tổng buông một câu nửa đùa nửa thật.

Bất chợt, từ sau chiếc ghế bành, hai cái đầu thò lên một cách đột ngột như thể đã ẩn nấp từ lâu chờ khoảnh khắc hoàn hảo để xuất hiện. Không ai khác, đó là Watame và Nene, hai cái tên gắn liền với sự ngây thơ và ngờ nghệch. Nhưng lúc này, đôi mắt họ sáng rực như có ánh lửa thắp lên, tràn ngập sự ngưỡng mộ thuần khiết, khi nhìn thấy Ayame đang ngồi xổm, đầu đội mũ, miệng ngậm hoa, mắt che bởi kính đen.

\img{10-11c}

``Ng-Ngầu quá\dots'' cả hai đồng thanh, giọng run lên vì xúc động như thể vừa chứng kiến một vị anh hùng bước ra từ trong màn ảnh.

Kuon chẹp miệng, liếc mắt sang Subaru đứng kế bên, lẩm bẩm: ``Lại có thêm hai kẻ khờ dại nữa đây\dots''

Subaru chỉ thở dài, ánh mắt vừa bất lực vừa thương hại cho những người vừa ``sa lưới phong trào''.

Nhưng chưa dừng lại ở đó. Điều khiến cả hai người khóa trên thật sự sốc chính là hành động tiếp theo: Watame và Nene\dots\ quỳ sụp xuống trước mặt Ayame (hay đúng hơn là trước mặt ``chị đại giang hồ'' tự phong của họ).

``\textbf{Xin chị hãy thu nạp bọn em ạ!}'' cả hai cô gái đồng thanh, mặt nghiêm túc như đang xin gia nhập một tổ chức thần bí nào đó.

Nàng Vịt giật mình, ngơ ngác nhìn họ: ``Ơ\dots\ Ủa\dots?''

Cậu Tổng cúi xuống một chút để xem mình có nhìn lầm không, rồi nhăn mặt hỏi: ``\dots Hai đứa bị sao vậy?''

Ayame không hề tỏ vẻ bất ngờ. Trái lại, cô chậm rãi đứng dậy, khoanh tay trước ngực, phun ra câu đáp lạnh như băng: ``Nghe cho rõ đây. Ta là một con sói đơn độc. Băng đảng không phải gu của ta.'' Cô quay lưng ngoảnh mặt đi, động tác vừa khinh đời vừa\dots\ thiếu nghiêm túc đến mức khó tin.

Thế nhưng, điều đó chỉ càng khiến Watame và Nene cảm thấy cô ``ngầu lòi'' hơn. Hai cô gái lần lượt ngẩng đầu lên, thốt ra những lời thề độc như muốn chứng minh sự thành tâm không thể lay chuyển của mình.

``Làm ơn ạ!'' nàng Cừu van vỉ.

``Chúng em muốn gọi chị là `đại ca' ạ!'' nàng ``Hải cẩu'' tiếp lời.

\img{10-11d}

Kuon thở dài, giọng bất lực: ``\dots Dù không muốn nhưng anh phải đồng ý. Chẳng có giang hồ nào anh biết lại lẻ loi đi một mình đâu.''

Subaru nhăn mặt, ngạc nhiên khi cậu dường như theo phe của Ayame ``Đừng có tham gia vào mấy chuyện này chứ!''

``Anh chỉ đang nói lên sự thật thôi mà,'' cậu nhún vai. ``Em có thấy ai làm xã hội đen mà là `sói cô độc' không?''

Cô nàng ậm ừ, vẫn không thẻ chấp nhận thực tế phi lý trước mắt: ``Nhưng mà kể cả thế thì\dots!''

Nàng Quỷ sừng khựng lại trước những tiếng gọi ``đại ca'' đầy chân thành cùng ánh mắt tròn xoe chờ đợi của Watame và Nene. Thêm vào đó là một chút lý lẽ của Kuon. Cuối cùng, nàng Quỷ sừng cũng nhún vai, quay lại, nói bằng giọng thoải mái: ``Được thôi, sao cũng được. Hắn ta (Kuon) nói cũng có lý.''

Nàng Cừu và nàng ``Hải cẩu'' như vỡ òa trong sung sướng, bật dậy, hai tay giơ lên trời, hô vang: ``\textbf{Đại ca vạn tuế!!}''

Lập tức, Ayame giơ tay ra hiệu ngắt lời họ: ``Với một điều kiện.''

Cả hai nín thở.

``Nếu các ngươi muốn làm đồng bọn với ta\dots'' cô rút ra từ đâu đó hai chiếc mũ lưỡi trai cùng mẫu với chiếc cô đang đội, ``\dots thì các ngươi phải ăn mặc cho chỉnh tề vào.''

Nene và Watame hớn hở cầm lấy mũ, đội lên đầu một cách đầy tự hào, như thể mình vừa gia nhập hàng ngũ tinh anh.

Ayame thậm chí còn quay sang phía Kuon, chìa ra một chiếc mũ thứ ba: ``Ngươi thì sao, Cửu Viễn? Có muốn nhập hội với bọn này không?''

Cậu con trai nhìn cái mũ một lúc, rồi lắc đầu nhè nhẹ: ``Thôi, không cần đâu. Làm người tốt còn chưa ổn, nói gì đến làm người xấu.''

Cô nàng nhún vai, tặc lưỡi: ``Tùy ngươi thôi. Ta vẫn sẽ luôn đón chào ngươi.''

Subaru đứng bên cạnh, thở phào nhẹ nhõm: ``May quá, anh ấy không bị sa ngã\dots''

Từ hôm đó, một ``tổ chức'' mới thành lập trong văn phòng \hll, một thứ nửa vời giữa nhập vai và tâm thần nhẹ.

Ayame, trong vai ``đại ca'', được Kuon đặt biệt danh là Bách Quỷ Đại Đế.

Nene thì tự phong danh hiệu oai vệ không kém: Tiểu Linh Đào Tể, tay cầm chiếc cờ cam với vẻ dũng mãnh.

Còn Watame hùng hồn tuyên bố tên giang hồ của mình là Dương Đầu Mã Điện, với một lá cờ vàng chữ đen nắm trên tay, viết to tướng như vậy.

\img{10-11e}

Mặc dù đã biết rằng cậu Tổng vẫn còn tỉnh táo, nàng Tomboy vẫn không khỏi thở dài khi chứng kiến ba người kia cứ cư xử ngày càng kỳ dị. ``Mấy người mất trí hết rồi à?'' cô nói không môt chút kiêng dè.

Nhưng, chẳng ai thèm để ý tới lời nàng Vịt. Họ đang bận nói chuyện ``ngầu lòi'' với nhau như những bằng hữu trong giới giang hồ.

``Các ngươi trông khá là ngầu căng đét đấy.'' Ayame gật gù nhận xét.

Nene và Watame đồng thanh, cúi rạp người xuống: ``Cảm ơn đại ca ạ!!''

Nàng Oni chậm rãi quay sang: ``Ngươi thì sao, Cửu Viễn?''

Kuon khoanh tay, đáp ngắn gọn: ``Anh không có ý kiến gì cả.''

\ct

Việc xây dựng đội hình đã hoàn tất, giờ đến phần kế tiếp: huấn luyện phong thái và thần thái.``Bách Quỷ Đại Đế'' đứng trước hai ``đệ tử'' của mình, tay chống nạnh, khẽ hất cằm ra hiệu.

``Hai ngươi thử nhìn đểu cho ta xem nào.''

Cô làm mẫu trước. Chiếc kính râm kiểu cũ che gần nửa khuôn mặt, đôi mày nhíu lại, đôi môi mím chặt. Đặc biệt, bông hoa ngậm nơi miệng khiến cho cô có phần giống một tên xã hội đen\dots\ lãng mạn.

\img{10-11f1}

Kuon đứng sau lưng phì cười, chẹp miệng nhận xét: ``Trông giống như một bệnh nhân tâm thần hơn\dots''

Dù vậy, Nene và Watame lại không nghĩ thế. Họ vỗ tay bôm bốp như khán giả vừa chứng kiến một màn biểu diễn đỉnh cao, mắt lấp lánh ánh sao.

``Các ngươi cũng thử xem nào,'' nàng Quỷ sừng ra hiệu.

Không chần chừ, cả hai bước ra phía trước, dàn đội hình, hít một hơi thật sâu, rồi cùng ngẩng đầu lên, mắt đảo qua đối phương, miệng buông ra một tiếng ``hả'' đanh đá, như thể đang thách thức nhau đánh một trận sống mái.

\img{10-11f2}

Kuon nghiêng đầu, bình thản nói: ``Anh nhìn hai đứa như đang chuẩn bị đánh nhau rồi đấy.''

Subaru đứng bên cạnh đáp lại: ``Chả phải mục đích nhìn đểu là như vậy sao?''

Cả ba người trong nhóm ``giang hồ'' cứ tiếp tục lườm nhau thêm một chút nữa, khiến cho nàng Vịt bị thu hút. Đôi môi nở một nụ cười nhỏ: ``Ơ kìa, trông vui thế nhỉ\dots\ Em có thử được không?''

``Cứ thoải mái đi. Có gì anh bảo kê,'' cậu Tổng đáp tỉnh bơ.

Được đà, Subaru tiến lại gần ba người đang tập luyện. Cô dừng lại trước mặt họ, hít một hơi thật sâu, cố nhướng cằm lên, cau mày lại\dots\ rồi phóng ánh mắt ``đầy sát khí'' vào đối phương.

Nhưng khuôn mặt của cô trong con mắt của ba cô gái lại có đôi phần\dots\ kỳ quặc.

Gương mặt của nàng Vịt chẳng gợi chút đáng sợ nào. Mắt cô mở to, ánh sáng phản chiếu rõ cả logo như hoạt hình trong con ngươi, gò má ửng hồng, miệng hé ra trông như đang hỏi ``cái shuba gì cơ?''. Nét mặt của cô không ra dáng một tên giang hồ, mà lại giống hệt một cô bé vừa bị mắng oan, hoang mang cực độ giữa bầu không khí căng thẳng.

\img{10-11f3}

Nene, Watame và Ayame dừng lại, giương cặp mắt có phần nửa thương nửa hài nhìn cô. Kuon ở đằng sau thở dài, phán ra một câu như đấm thẳng vào lòng tự trọng của cô nàng: ``Trông em chẳng khác gì một con nghiện hút xì ke ấy.''

Và thế là hết. Subaru bật khóc tức tưởi, gào lên: ``\textbf{Sao mấy người lại cứ xua đuổi Subaru ra thế hả!?}'' và chạy vụt ra ngoài như một cơn gió mùa hạ, để lại sau lưng là ánh mắt ngẩn ngơ của cả nhóm.

``\dots Chị ấy bị sao vậy, Kuon-senpai?'' Nene hỏi khẽ.

``Tại mấy đứa trưng cái mắt cá ươn ra nên em ấy sợ đấy,'' Kuon đáp lại, giọng nửa châm biếm nửa thật tình.

``Chứ không phải là do câu nói thẳng tính của ngươi sao?'' Ayame nhe răng.

``\dots Được rồi, một phần cũng vì cái đó nữa.''

``Subaru-senpai sẽ ổn chứ?'' Watame hướng nhìn về cánh cửa, tỏ vẻ lo lắng.

``Em ấy sẽ ổn thôi.''

Nàng Quỷ sừng hất tóc, vờ như chẳng quan tâm tới nàng Vịt vừa khóc lóc ấy: ``Quên cái ả đó đi. Mấy đứa làm cũng ra trò phết.''

Hai đệ tử cảm ơn ``chị đại'' rối rít. Kuon thì không nói gì thêm, chỉ khoanh tay đứng nhìn.

Ayame quay lại, mắt lóe lên vẻ háo hức: ``Giờ thì tới bước cuối\dots\ \textbf{Nẹt pô nào mấy đứa!}''

``Ờ thôi, anh miễn nha,'' Kuon giơ tay từ chối. ``Anh không gấu đến mức đó đâu.''

``Ngươi thích thì ngươi có thể không tham gia được mà,'' nàng Quỷ sừng trả lời, nụ cười khẩy hiện rõ trên môi.

Watame và Nene thì mắt sáng rực như trẻ con lần đầu thấy pháo hoa: ``Nẹt pô thật hả chị?''

``Miễn là đừng làm gì quá điên rồ nha là được,'' cậu nhắc khéo, giọng hơi lo.

\ct

Dưới ánh hoàng hôn vào những ngày cuối cùng của mùa thu đổ dài trên mặt biển Điện tử, những cơn gió biển thổi vào mang theo hương mặn mòi dịu dàng của đại dương, càng khiến cho mọi người cảm tưởng rằng một mùa đông nữa đã sắp tới.

Trên đoạn đường vắng bên bờ, ``nhóm giang hồ'' đang\dots\ vi vu bằng những chiếc ``xe máy giả''. Cụ thể hơn, chỉ có Kuon là đi xe thật, một chiếc xe đạp mà cậu thuê được. Còn ba người kia\dots\ thì đẩy khung xe máy tự chế, chân chạy dưới đất. Động cơ duy nhất của họ là bằng sức lực dung nạp từ cơm và lòng nhiệt huyết.

``Ôi, làn gió mát chưa kìa,'' Ayame hít một hơi đầy khoan khoái.

Kuon chậm rãi gật đầu: ``Cũng sắp đến đông rồi nhỉ.''

Đằng sau họ, Nene và Watame đang chạy hụt hơi, chân như rã rời vì kéo ``xe máy'' bằng chân trần.

\img{10-11g}

Nàng Quỷ sừng quay đầu lại, gọi vọng: ``\textbf{Mấy ngươi có thấy phê không nào?}''

``C-Có ạ\dots'' nàng Cừu vừa thở dốc vừa nói.

``T-Tuyệt lắm\dots'' nàng Gấu đáp lại bằng một nụ cười méo xệch, từng nhịp thở ngắt quãng.

Cậu Tổng liếc nhìn. liền phì cười: ``Mát mà sao mặt nghệt thế kia?''

``K\dots Không có gì ạ\dots'' cả hai cùng trả lời, cố giấu vẻ kiệt sức.

Ayame cười lớn, giơ tay lên: ``\textbf{Được! Lên đi! Cùng đi qua đường chân trời nào!}''

Watame và Nene không biết lấy đâu ra sức nữa, lập tức hét lên một tiếng đồng lòng như nhận hiệu lệnh từ đại ca.

Và thế là, dưới ánh chiều rực vàng cuối tháng Mười Một, bốn con người - một người đạp xe, ba người ``nẹt pô'' bằng chân - chạy dọc ven biển, tiếng cười của họ vang lên hòa cùng sóng vỗ. Mỗi bước chạy đều nhẹ nhàng nhưng đầy ý nghĩa, như thể họ đang sống trọn vẹn một thời khắc điên rồ nhất của tuổi trẻ.

Họ cứ thế chạy mãi, chạy đến khi hình bóng của họ nhỏ dần, mờ dần\dots\ rồi biến mất nơi đường chân trời xa xăm, nơi ánh hoàng hôn cuối thu dần buông xuống thành phố.

%==========================================TRUYỆN PHỤ=========================================
\omk

\dots Đó ít nhất là điều mà nhóm ``giang hồ tự phong'' ấy muốn tin là mình đang làm.

Nhưng thực tế thì không hào nhoáng như phim ảnh. Giữa chiều muộn, bên bờ biển thành phố Điện tử, ba cô gái đang thở hổn hển, mồ hôi đẫm trán, chân mỏi nhừ như vừa chạy hết mười vòng sân vận động.

Chỉ có một người vẫn thong dong đạp xe là Kuon, trong bộ đồ đi làm thường nhật và thái độ bình thản như thể đang đi chợ.

``Mệt chưa?'' cậu cười nhẹ, nhìn ba người đang loạng choạng phía sau.

``C-Chưa\dots \! Còn lâu!'' Ayame gắt lên, giọng lạc hẳn vì hụt hơi.

Nene và Watame nối tiếp trong tiếng thở dốc: ``Nene\dots\ sắp chạy hết nổi rồi\dots'' / ``Em\dots\ cũng thế\dots''

Nàng Quỷ sừng vẫn cố ra vẻ, dù lưỡi của cô đã bắt đầu thè ra vì mệt: ``Các ngươi\dots\ yếu quá\dots! Nhìn gương ta\dots\ mà xem đây!''

Chưa kịp thể hiện thêm, Kuon chợt chỉ tay về phía trước, nói to: ``À, Suba-chan kìa!''

Thấy tín hiệu từ cậu, nàng Tomboy liền tiến lại gần với một túi nước mát lạnh. Trông thấy ba ``giang hồ'' như sắp gục đến nơi, cô không nói gì, chỉ đưa nước ra. Như bản năng sinh tồn, cả ba giật lấy ngay, uống ừng ực như vừa từ sa mạc về, không quên nói ``Cảm ơn chị!'' giữa những tiếng thở phì phò.

\img{10-11h}

``Mấy đứa vẫn còn non lắm,'' Subaru nói nhẹ nhàng, nhưng đủ khiến ba người im lặng.

Kuon dựng chiếc xe đạp, rồi khẽ khàng thở dài, tiếp lời cô: ``Thực sự mà nói\dots\ Làm giang hồ còn khó hơn thế này nhiều.''

Cả Nene và Watame đều mở to mắt, giọng ngạc nhiên: ``Thế ạ? Vậy\dots\ theo anh, giang hồ sẽ làm những gì?'' / ``Nene cũng muốn biết xem họ thực ra như thế nào!''

``Cái này anh có thể liệt kê đến già. Nhưng để ngắn gọn ấy, mấy ý này thôi.''

``Cụ thể là gì?'' nàng Quỷ sừng chồm tới.

``Một,'' cậu giơ ngón tay lên `` ăn mặc sành điệu. Và khi anh nói `sành điệu' thì ý anh nói không khác gì một lũ bặm trợn, mặt mũi lúc nào cũng hằm hằm ấy, tóc tai thì xù lên xù xuống, nhuộm màu lòe loẹt ấy.''

``Giống như mọi người bây giờ thôi,'' Subaru phụ họa.

``Chả phải\dots\ em cũng đã làm thế sao?'' Watame nhìn lại từ đầu tới chân, nghiêng đầu thắc mắc.

``Nếu chỉ bề ngoài thôi là chưa đủ,'' Kuon đáp. ``Mấy đứa biết giang hồ thì phải làm gì không?''

``Là mấy thứ mà ta không nên làm đấy,'' nàng Vịt chen vào.

``Ví dụ như\dots?'' Ayame hỏi, nhưng trong lòng đã bắt đầu thấy hơi bất an.

``Cho vay nặng lãi, đòi nợ thuê\dots'' Kuon mở đầu, xòe từng ngón tay.

``Eep\dots'' Nene siết chặt tay nắm mũ.

``Cờ bạc trá hình, buôn ma túy\dots'' Subaru nối tiếp, giọng đều đều.

``Oái\dots'' Watame rụt cổ, rõ ràng bị sốc.

``Can dự vào nội bộ công ty, bảo kê cho tội phạm, đánh nhau với các băng khác, tranh giành lãnh địa\dots'' Kuon liệt kê thêm, không chút ngập ngừng.

``Thậm chí còn xử nhau trong nội bộ nếu có mâu thuẫn lợi ích,'' Subaru kết lại, giọng nghiêm túc hiếm thấy.

Ayame đứng như hóa đá, bông hoa ngậm nơi miệng suýt thì rơi xuống. ``Ối chà\dots\ Vậy những gì mà tôi đọc trong cuốn sách đó\dots''

``\dots chỉlà những thứ được lý tưởng hóa thôi. Truyện tranh mà. Làm gì có thật.'' Subaru nói, giọng nhẹ nhưng không kém phần dứt khoát. ``\textit{Trừ khi bà tính làm thật ở trong môi trường này\dots}'' cô nàng nghĩ thầm, hoàn toàn không muốn động viên cô đi theo con đường này.

``Anh cảm thấy em dễ bị ảnh hưởng bởi mấy thứ này lắm đấy, Ayame-chan ạ,'' Kuon thở dài, rồi nhìn vào mắt cô. ``Cũng chẳng trách được, em vẫn còn làm một cô oni hiếu kỳ, muốn tìm hiểu thêm về thế giới con người bọn nah mà. Nhưng mà theo anh ấy\dots''

Cậu và Subaru đồng thanh, giọng nhẹ nhàng hẳn như những người đang đưa họ về đúng con đường thiện chi: ``Ba đứ anhư là chính mình đi là tốt nhất.''

Không cần giải thích nhiều, cậu Tổng tiến tới, nhẹ nhàng tháo mũ và kính râm trên mặt Ayame.
``Hãy cứ là con người lương thiện\dots'' cậu thì thầm.

Nàng Tomboy cũng bước tới, làm điều tương tự với Nene và Watame, tay vỗ nhẹ lên vai họ: ``\dots Vì chỉ như thế thì mới có thể trong sạch được. Chúng ta là idol mà, không phải tội phạm đâu.''

Ba người con gái, vốn vừa nãy còn hung hăng, giờ như những học sinh bị thầy cô nhắc nhở đúng lúc, giờ đã cúi đầu. Ayame thốt lên: ``Ta\dots\ Ta hứa sẽ trở thành một con người có ích, bang chủ!''

``\dots Ờm, anh không phải là bang chủ\dots'' Kuon gãi đầu. ``Gọi anh bằng tên như mọi khi là được rồi.''

``Cứ xưng hô như bình thường là được rồi,'' Subaru mỉm cười, ánh mắt dịu dàng.

Nene và Watame cũng cúi đầu, đồng thanh: ``Xin hai người tiếp tục chỉ bảo chúng em ạ!''

``Rất hân hạnh.''

Và từ khoảnh khắc ấy, Ayame, Nene, Watame đã dần hiểu rằng làm một người bình thường, sống tử tế, giúp đỡ người khác và không gây rắc rối\dots\ vẫn là con đường được nhiều người quý trọng hơn cả.

Không cần đóng vai ``đại ca'', không cần ``nẹt pô'' hay đeo kính râm tối sầm, chỉ cần sống đúng với bản thân, trọn vẹn và chân thành, thế là đủ.

Vì cuối cùng, họ cũng nhận ra, như câu thơ của Tố Hữu viết trong bài ``Từ ấy'':

\tabs\tabs\tabs\tabs\tabs  ``Từ ấy trong tôi bừng nắng hạ\\
\tabs\tabs\tabs\tabs\tabs  Mặt trời chân lý chói qua tim.''

Và đúng vậy, mặt trời hôm đó đang dần lặn ở đường chân trời xa xa, để lại phía sau là một hoàng hôn ấm áp, thứ ánh sáng mà cả nhóm giang hồ nửa mùa đã tìm được, không phải từ quyền lực hay sợ hãi, mà từ những lời nhắc nhở tử tế và chân thành.

\end{document}